% 재난 릴레이 드론 배송을 위한 계층형 다중 에이전트 강화학습 (한글본)
% 대상 버전: v3 (26_08_31_19)
% 컴파일: xelatex 또는 lualatex (kotex 사용)
\documentclass[12pt]{article}
\usepackage{kotex}
\usepackage[margin=1in]{geometry}
\usepackage{setspace}
\usepackage{amsmath, amssymb}
\usepackage{booktabs}
\usepackage{graphicx}
\usepackage{url}
\doublespacing
\pagestyle{empty}

\title{재난 상황의 릴레이 드론 네트워크 기반 물자 배송을 위한\\계층형 다중 에이전트 강화학습}
\author{고광종 \and 최현석 \and 정태수\thanks{교신저자}\\
고려대학교 산업경영공학과}
\date{}

\begin{document}
\maketitle

\begin{abstract}
재난으로 통신 인프라가 파괴된 지역에서는 드론의 직접 통신 반경이 배송 가능 영역을 제한한다.
각 드론이 서로의 통신을 중계하는 릴레이 드론 네트워크는 이 제한을 완화하지만, 중계 의무와
배송 처리량이 서로 경쟁하는 새로운 운영 문제를 만든다. 본 연구는 동질적 드론 편대가 물자를
배송하면서 동시에 다중 홉 중계 사슬을 형성하여 모든 기체가 제어 센터와 항시 연결되도록 하는
문제를 다룬다. 목적함수는 전 목적지 배송 완료 시간(makespan)의 최소화다.

연결성을 보상 페널티로 다루는 기존 방식은 회피 불가능한 페널티를 학습 신호에 주입한다는 것이
관찰되어, 본 연구에서는 연결을 유지하는 최대 변위까지 이동을 축소하는 최소 섭동 투영으로
연결성을 환경의 하드 제약으로 강제하였다. 또한 스텝 단위 연속 제어만으로는 적재량 소진 후
거점으로 복귀하여 재적재하는 왕복 사이클이 학습되지 않는 현상이 확인되어, 상위가 배송지 선택과
복귀 및 중계 대기를 이산적으로 결정하고 하위가 2차원 추력을 담당하는 계층 구조를 제안한다.

통신 반경 300, 목적지 50곳, 드론 4대 설정의 홀드아웃 30개 인스턴스 $\times$ 3회 반복 실험에서
제안 방법은 배송 23.97건(표준오차 0.67)으로 탐욕 휴리스틱 23.04건(표준오차 1.01)과 통계적으로
구분되지 않았다. 반면 인스턴스 간 표준편차는 9.62에서 6.34로 감소하였고, 제어 센터로부터의
최대 도달 거리는 644에서 778로 증가하였으며, 연결 제약에 의한 이동 취소는 3,338회에서 1,373회로
59\% 감소하였다. 즉 제안 방법의 이점은 평균 처리량이 아니라 성능의 일관성과 중계 네트워크
활용도에 있다. 통신 단절은 학습 및 평가 전 구간에서 한 건도 발생하지 않았다.
\end{abstract}

\section{서론}

재난 상황에서 긴급 물자를 수송할 때 지상 운송 수단은 인프라 파괴로 사용이 어렵다. 드론 배송이
대안으로 검토되어 왔으나, 재난 지역에서는 통신 인프라 역시 파괴되어 있다는 점이 물류 문헌에서
충분히 다루어지지 않았다. 긴급 물자는 배송 완료 중요도가 높아 배송 상황 추적이 필요하고,
상하역 시 세밀한 제어를 위해 영상 데이터의 실시간 전송이 요구된다. 따라서 임무 수행 중 통신
연결이 상시 유지되어야 한다.

각 드론이 다른 드론의 통신을 중계하는 릴레이 드론 네트워크는 편대의 임무 반경을 직접 통신
반경 너머로 확장한다. 그러나 이 구조는 운영상의 상충을 만든다. 중계 위치를 지키는 드론은
배송을 수행하지 못하므로, 중계에 몇 대를 배정할지가 곧 배송 처리량과의 교환이 된다. 본 연구의
사전 실험에서 이 교환의 최적점이 인스턴스마다 달라지는 구간이 존재함이 확인되었으며(제5.1절),
이 구간이 고정 규칙으로는 대응할 수 없고 학습된 적응적 배분이 기여할 수 있는 영역이다.

본 연구의 기여는 세 가지다. 첫째, 통신 반경을 축으로 문제 난이도를 훑어 규칙 기반 정책이
실패하기 시작하는 구간과 고정 역할 배분이 원리적으로 최적일 수 없는 구간을 특정하였다.
둘째, 연결성을 보상이 아니라 최소 섭동 투영으로 강제하여 통신 단절을 구조적으로 제거하였다.
셋째, 배송지 선택과 저수준 기동을 분리한 계층형 다중 에이전트 강화학습 프레임워크를 제안하고,
탐욕 휴리스틱 대비 배송량은 동등하나 성능 분산이 34\% 낮고 도달 범위가 21\% 넓음을 보였다.

\section{선행연구}

\subsection{드론 배송 라우팅}

Murray와 Chu(2015)는 트럭 1대와 드론 1대가 협업하는 flying sidekick TSP를 정식화하고 makespan
최소화를 목적으로 하는 혼합정수계획 모형과 휴리스틱을 제안하였다. Murray와 Raj(2020)는 이를
이종 드론 편대로 확장하였고, Agatz 등(2018)은 동일 노드에서의 발사와 회수를 허용하는 정식화와
정확 동적계획 알고리즘을 제시하였다. Dorling 등(2017)은 드론 전용 다중 왕복 차량경로문제를
정의하고 적재 중량과 배터리 중량에 선형인 에너지 모형을 도입하였다. 이 계열의 연구는 통신을
비용 없이 항상 가능한 것으로 가정한다.

\subsection{UAV 중계 네트워크}

Zeng 등(2016)은 UAV를 이동 중계 노드로 활용하는 mobile relaying을 제안하고 전력 할당과 궤적을
교번 최적화하였으나 단일 중계 UAV를 가정하였다. Yanmaz(2022)는 임무 UAV와 중계 UAV의 역할을
분리하고 Steiner 점 기반으로 최소 수의 중계 UAV 배치를 계산하는 2단계 알고리즘을 제안하였다.
Ding 등(2022)은 다중 홉 UAV 중계망에서 궤적, 주파수 할당, 다음 홉 선택을 QMIX 기반 다중 에이전트
강화학습으로 동시에 결정하였다. 이 계열의 목적함수는 처리량과 커버리지이며, 적재량 제약이나
재적재 왕복과 같은 물류 요소는 다루지 않는다.

\subsection{연구 공백}

임무 수행과 통신 중계 사이의 역할 전환 자체는 다중 로봇 탐사 문헌에서 확립된 개념이다.
Cesare 등(2015)은 explore, meet, sacrifice, relay의 네 상태를 정의하고 relay 역할 로봇이 정지하여
중계 노드로 기능하게 하였다. 따라서 역할 전환을 본 연구의 신규성으로 주장할 수 없다.
저자들이 조사한 범위에서 확인되지 않은 것은, 이 역할 전환을 \emph{용량 제약이 있고 재적재를 위한
반복 왕복이 필요하며 makespan을 최소화하는} 배송 문제로 옮긴 사례다. 이 설정에서는 중계 의무가
적재 처리량과 직접 경쟁하며, 편대가 보급을 위해 거점으로 반복 복귀해야 한다.

\subsection{학습 방법론}

Huang과 Ontañón(2022)은 실행 불가능한 행동을 보상 페널티로 억제하는 방식과 마스킹으로 차단하는
방식을 통제 비교하여, 문제 규모가 커질수록 전자의 에피소드 리턴이 붕괴함을 보였다. Dalal 등
(2018)은 연속 행동에서 제약을 만족하는 최소 섭동을 해석적으로 구하는 안전 계층을 제안하였다.
Ng 등(1999)은 포텐셜 기반 보상 쉐이핑이 최적 정책을 보존하는 필요충분 형태임을 증명하였고,
그렇지 않은 스텝당 양의 보상이 목표 완수 대신 순환 행동을 유발하는 사례를 제시하였다.
Zhou 등(2019)은 회전을 각도 스칼라로 표현하면 신경망 학습에 불리한 불연속이 발생하며 SO(2)의
연속 표현이 $[\cos\theta, \sin\theta]$임을 보였다. 본 연구의 설계는 이들 결과를 반영한다.

\section{문제 정의}

제어 센터 1곳과 목적지 $|C|$곳이 $1000 \times 1000$ 평면에 주어진다. 동질적 드론 $|K|$대는
제어 센터에서 적재량 $Q$만큼 물자를 싣고 출발한다. 드론 $i$가 목적지 $j$의 상호작용 반경
$r$ 이내에 진입하면 배송이 이루어지고 적재량이 1 감소하며 $S$스텝의 하역 시간이 소요된다.
적재량이 0인 드론이 제어 센터에 복귀하면 적재량이 $Q$로 회복되고 동일하게 $S$스텝이 소요된다.

통신 그래프 $G_t = (V, E_t)$는 매 시점 $t$의 드론 위치로 정의된다. 두 노드는 거리가 통신 반경
$R$ 이내일 때 연결되며, 제어 센터도 노드에 포함된다. 제약은 다음과 같다.
\begin{equation}
\forall t, \ \forall i \in K: \ i \text{는 } G_t \text{에서 제어 센터와 연결되어야 한다.}
\end{equation}
목적함수는 모든 목적지의 배송이 완료되는 시점 $T^*$의 최소화다.

수리계획 정식화의 어려움을 밝혀 둔다. 노드 기반 라우팅 정식화에서 이동 중인 드론은 좌표를
갖지 않으므로 매 시점 연결성을 판정할 수 없다. 이를 메우려면 공간을 격자화하거나 호버링 후보
지점을 사전 지정해야 하며, 후자는 그 자체가 배치(placement) 문제다. 동일 드론이 배송과 중계를
겸하므로 Yanmaz(2022)처럼 배치를 하위 문제로 분리할 수도 없어 배치와 경로가 결합된 문제가 된다.

\section{방법론}

\subsection{연결성의 하드 제약화}

선행 구현에서 연결성은 단절 시 대량의 페널티를 부여하고 에피소드를 강제 종료하는 방식으로
다루어졌다. 이 방식의 결함이 본 연구에서 확인되었다. 하역 중이어서 이동 의사가 없던 드론이
동료의 이탈로 고립되는 경우, 해당 드론은 어떤 행동으로도 회피할 수 없는 페널티를 받는다.
회피 불가능한 페널티는 학습 신호를 오염시킨다.

본 연구는 연결성을 환경의 하드 제약으로 전환하였다. 제안된 변위 $\Delta$에 대해 배율
$s \in [0,1]^{|K|}$를 정하여 최종 위치를 $x + \Delta \odot s$로 두되, 결과 그래프에서 전원이
연결되도록 한다. 초기값은 $s = \mathbf{1}$이며, 연결이 성립하지 않는 동안 제한 시 연결 드론 수가
가장 많아지는 드론을 선택하여 그 배율을 이분 탐색으로 얻은 최대 안전값까지 축소한다. 동점은
무작위로 깬다. 결정론적으로 깰 경우 중계 사슬을 고정하는 드론이 반복 선택되어 불이익이
집중되는 현상이 관찰되었기 때문이다. 모든 배율을 0으로 되돌리면 직전 위치가 되고 이는 연결
상태이므로, 절차는 항상 연결된 구성에서 종료한다.

\subsection{계층 구조}

스텝 단위 연속 제어만으로는 적재량 소진 후 복귀하여 재적재하는 왕복 사이클이 학습되지 않았다
(제5.3절). 이 사이클은 수백 스텝에 걸친 시간 지평을 갖는다. 본 연구는 의사결정을 두 층으로
분리한다.

\textbf{상위 정책.} $H$스텝마다 또는 현재 하위 목표가 무효화될 때 각 드론에 대해 $K+2$개 중
하나를 선택한다. 선택지는 자기 기준 최근접 미배송지 $K$개, 제어 센터 복귀, 현 위치 유지다.
목적지 전체를 행동으로 두면 행동 공간이 인스턴스 크기에 종속되므로 최근접 $K$개로 고정한다.
현 위치 유지를 명시적 행동으로 두는 것이 설계의 핵심이다. 중계 역할은 본질적으로 이동하지
않는 것이며, 이를 명시적 선택지로 두어야 상위가 역할 분담을 직접 결정할 수 있다.
보상은 팀 목적함수(배송, 재적재, 시간 페널티)만 받는다.

\textbf{하위 정책.} 매 스텝 2차원 추력 $(v_x, v_y) \in [-1,1]^2$를 출력한다. 노름이 1을 넘으면
정규화하며 변위는 $v \cdot v_{\max}$다. 각도 스칼라 표현은 SO(2)의 불연속 표현이어서 물리적으로
동일한 두 헤딩이 행동 공간의 양 끝에 놓이고, 최대 속도에 도달하려면 tanh 포화가 필요하여
최대 엔트로피 목적함수가 이를 억제한다. 본 연구의 계측에서 각도 표현 정책의 평균 명령 속도는
최대치의 0.511에 그쳤다. 보상은 하위 목표까지의 거리 감소분, 도달 보너스, 연결 보정으로
잘려나간 변위에 비례하는 페널티로 구성되며 팀 보상은 받지 않는다.

두 층 모두 중앙집중 학습과 분산 실행(CTDE) 구조의 soft actor-critic으로 학습된다. 상위는
범주형 정책의 이산 SAC, 하위는 tanh 가우시안 정책의 연속 SAC를 사용한다.

\section{실험}

\subsection{문제 난이도 축}

통신 반경을 축으로 탐욕 정책의 성능을 측정하였다(표~\ref{tab:sweep}). 마지막 열은 최선의
\emph{고정} 중계 대수와, 인스턴스마다 최적 대수를 선택할 수 있다고 가정했을 때의 격차다.

\begin{table}[h]
\centering
\caption{통신 반경별 탐욕 정책 성능 (인스턴스 12개, 배송 건수/50)}
\label{tab:sweep}
\begin{tabular}{rrrrrrrr}
\toprule
반경 & 중계 0대 & 1대 & 2대 & 3대 & 최선 고정 & 인스턴스별 최적 & 격차 \\
\midrule
500 & 49.3 & 50.0 & 50.0 & 36.4 & 50.0 & 50.0 & $+0.0$ \\
400 & 40.6 & 48.4 & 43.9 & 17.9 & 48.4 & 49.4 & $+1.0$ \\
\textbf{300} & 21.0 & 22.7 & 18.5 & 9.3 & 22.7 & \textbf{28.2} & $\mathbf{+5.6}$ \\
250 & 14.4 & 15.2 & 12.6 & 7.8 & 15.2 & 19.1 & $+3.8$ \\
200 & 10.8 & 9.7 & 8.2 & 5.9 & 10.8 & 11.8 & $+1.0$ \\
150 & 7.0 & 5.8 & 4.6 & 4.2 & 7.0 & 7.2 & $+0.2$ \\
\bottomrule
\end{tabular}
\end{table}

반경 400까지는 탐욕이 문제를 사실상 해결하므로 학습 기반 방법이 기여할 여지가 없다. 반경
300에서 성능이 절반 이하로 떨어지는 동시에 격차가 최대가 된다. 이는 고정 중계 대수로는 대응할
수 없고 인스턴스마다 다른 배분이 필요함을 의미한다. 반경 200 이하에서는 격차가 다시 줄어드는데,
어떤 배분으로도 성능을 내기 어려운 영역이기 때문이다. 이후 실험은 반경 300에서 수행하였다.

\subsection{주 결과}

홀드아웃 인스턴스 30개를 3회씩 굴려 총 90개 롤아웃으로 평가하였다(표~\ref{tab:main}).
반복은 연결성 투영의 무작위 동점 처리에서 오는 분산을 평균하기 위한 것이다.

\begin{table}[h]
\centering
\caption{주 결과 (홀드아웃 30 인스턴스 $\times$ 3회, 통신 반경 300)}
\label{tab:main}
\begin{tabular}{lrrrrrr}
\toprule
방법 & 배송 & 표준편차 & 표준오차 & 재적재 & 최대 도달 & 이동 취소 \\
\midrule
무작위 정책 & 3.7 & 2.5 & 0.26 & 0.0 & 372 & 17 \\
탐욕 & 23.04 & 9.62 & 1.01 & 2.12 & 644 & 3{,}338 \\
탐욕 + 중계 1대 고정 & 14.24 & 5.29 & 0.56 & 0.96 & 524 & 2{,}652 \\
규칙 상위 + 학습 하위 & 23.09 & 9.25 & 0.97 & 2.50 & 695 & 1{,}303 \\
\textbf{제안 (계층, $H=20$)} & \textbf{23.97} & \textbf{6.34} & \textbf{0.67} & 2.34 & \textbf{778} & \textbf{1{,}373} \\
제안 (계층, $H=10$) & 22.52 & 5.75 & 0.61 & 2.11 & 788 & 1{,}322 \\
\bottomrule
\end{tabular}
\end{table}

배송 건수에서 제안 방법은 23.97건으로 탐욕의 23.04건을 앞서지만, 두 값의 표준오차가 각각
0.67과 1.01이므로 차이 $0.93$은 통계적으로 유의하지 않다. \textbf{배송 처리량 관점에서 제안
방법은 탐욕 휴리스틱과 동등하다고 보고하는 것이 정확하다.}

유의한 차이는 다른 지표에서 나타난다. 인스턴스 간 표준편차가 9.62에서 6.34로 34\% 감소하여
성능의 일관성이 개선되었다. 제어 센터로부터의 최대 도달 거리는 644에서 778로 증가하였는데,
통신 반경이 300임을 고려하면 직접 통신 반경의 2.6배까지 중계 사슬이 뻗었음을 뜻한다. 연결
제약에 의한 이동 취소는 3,338회에서 1,373회로 59\% 감소하였다. 즉 제안 방법은 더 멀리 나아가면서
제약과는 덜 충돌한다.

통신 단절은 학습 전 구간과 평가 전 구간에서 한 건도 발생하지 않았다. 이는 제안한 하드 제약이
설계대로 작동함을 확인한다.

\subsection{계층 분해의 기여}

계층 분해가 겨냥한 병목은 재적재 왕복이었다. 단일 계층 구조에서 학습된 정책의 재적재 횟수는
에피소드당 0.0에서 0.4회에 머물렀다. 드론 4대에 적재량 5이므로 재적재 없이는 20건이 구조적
상한이며, 관측된 배송량은 그 아래에 갇혀 있었다. 계층 구조에서는 재적재가 2.34회로 탐욕의
2.12회와 대등한 수준으로 회복되었다.

상위 결정 주기 $H$에 대한 감도를 확인한 결과 $H=20$이 $H=10$보다 우세하였다(23.97 대 22.52).
재결정이 잦으면 하위 목표가 자주 바뀌어 하위 정책이 목표에 도달하기 전에 목표가 교체된다.

\subsection{절제 실험}

여러 설계 요소를 개별적으로 검증하였다. 중계 간격에 대한 포텐셜 기반 쉐이핑, 중계 필수 드론에
대한 팀 보상 가중 재분배, 엔트로피 계수 하한, 동료 전원 관측으로의 확장은 모두 배송 성능을
개선하지 못하였다. 뭉침 페널티는 제거하는 편이 우세하여 최종 설계에서 제외하였다.

행동 표현의 교체는 배송 성능을 바꾸지 않았으나 학습 안정성에 뚜렷한 효과가 있었다. 각도 표현
정책은 정점 이후 홀드아웃 성능이 40\% 하락한 반면 벡터 표현 정책은 8\% 하락에 그쳤고, 정점 도달
시점도 4배 늦어졌다.

\section{결론}

본 연구는 재난 상황의 릴레이 드론 배송을 계층형 다중 에이전트 강화학습으로 접근하였다.
연결성을 최소 섭동 투영으로 하드 제약화하여 통신 단절을 구조적으로 제거하였고, 배송지 선택과
저수준 기동을 분리하여 단일 계층 구조에서 학습되지 않던 재적재 왕복 사이클을 회복하였다.

결과는 제한적이다. 배송 처리량에서 제안 방법은 탐욕 휴리스틱과 통계적으로 구분되지 않는다.
학습 기반 방법이 잘 조정된 휴리스틱을 해의 품질로 능가하지 못하는 것은 신경 조합최적화 문헌에서
보고되어 온 바와 일치한다. 제안 방법의 이점은 성능 분산의 감소(34\%)와 중계 네트워크 활용도의
증가(도달 거리 21\% 증가, 이동 취소 59\% 감소)에 있으며, 이는 재난 대응이라는 응용 맥락에서
평균 처리량과 별개로 의미를 갖는다.

한계는 세 가지다. 첫째, 학습이 고정된 지도 하나에서 이루어져 인스턴스 암기의 여지가 있다.
평가는 학습에 사용하지 않은 인스턴스에서 수행하였으나, 에피소드마다 인스턴스를 무작위화하는
학습은 향후 과제로 남는다. 둘째, 어떤 방법도 50개 목적지 전량 배송을 달성하지 못하여 목적함수인
makespan이 정의되지 않았고, 배송 건수를 대리 지표로 사용하였다. 셋째, 통신을 반경 기반 원판
모형으로 단순화하였으며 지형에 의한 차폐를 고려하지 않았다.

향후 연구로는 두 계층의 동시 미세조정, 학습 시 인스턴스 무작위화, 그리고 실시간 재결정 능력을
정량화하기 위한 스텝당 추론 지연의 측정을 계획한다.

\section*{References}
\noindent Agatz, N., Bouman, P., \& Schmidt, M. (2018). Optimization approaches for the traveling
salesman problem with drone. \textit{Transportation Science, 52}(4), 965--981.

\noindent Cesare, K., Skeele, R., Yoo, S.-H., Zhang, Y., \& Hollinger, G. (2015). Multi-UAV
exploration with limited communication and battery. In \textit{Proceedings of the IEEE
International Conference on Robotics and Automation} (pp. 2230--2235).

\noindent Dalal, G., Dvijotham, K., Vecerik, M., Hester, T., Paduraru, C., \& Tassa, Y. (2018).
Safe exploration in continuous action spaces. \textit{arXiv preprint arXiv:1801.08757}.

\noindent Ding, R., Chen, J., Wu, W., Liu, J., Gao, F., \& Shen, X. (2022). Packet routing in
dynamic multi-hop UAV relay network: A multi-agent learning approach. \textit{IEEE Transactions
on Vehicular Technology, 71}(9), 10059--10072.

\noindent Dorling, K., Heinrichs, J., Messier, G. G., \& Magierowski, S. (2017). Vehicle routing
problems for drone delivery. \textit{IEEE Transactions on Systems, Man, and Cybernetics: Systems,
47}(1), 70--85.

\noindent Haarnoja, T., Zhou, A., Abbeel, P., \& Levine, S. (2018). Soft actor-critic: Off-policy
maximum entropy deep reinforcement learning with a stochastic actor. In \textit{Proceedings of the
35th International Conference on Machine Learning} (pp. 1861--1870).

\noindent Huang, S., \& Ontañón, S. (2022). A closer look at invalid action masking in policy
gradient algorithms. \textit{Proceedings of the International FLAIRS Conference, 35}.

\noindent Lowe, R., Wu, Y., Tamar, A., Harb, J., Abbeel, P., \& Mordatch, I. (2017). Multi-agent
actor-critic for mixed cooperative-competitive environments. In \textit{Advances in Neural
Information Processing Systems 30} (pp. 6379--6390).

\noindent Murray, C. C., \& Chu, A. G. (2015). The flying sidekick traveling salesman problem:
Optimization of drone-assisted parcel delivery. \textit{Transportation Research Part C: Emerging
Technologies, 54}, 86--109.

\noindent Murray, C. C., \& Raj, R. (2020). The multiple flying sidekicks traveling salesman
problem: Parcel delivery with multiple drones. \textit{Transportation Research Part C: Emerging
Technologies, 110}, 368--398.

\noindent Ng, A. Y., Harada, D., \& Russell, S. (1999). Policy invariance under reward
transformations: Theory and application to reward shaping. In \textit{Proceedings of the 16th
International Conference on Machine Learning} (pp. 278--287).

\noindent Park, J., Kwon, C., \& Park, J. (2023). Learn to solve the min-max multiple traveling
salesmen problem with reinforcement learning. In \textit{Proceedings of the 22nd International
Conference on Autonomous Agents and Multiagent Systems} (pp. 878--886).

\noindent Yanmaz, E. (2022). Positioning aerial relays to maintain connectivity during drone team
missions. \textit{Ad Hoc Networks, 128}, 102800.

\noindent Zeng, Y., Zhang, R., \& Lim, T. J. (2016). Throughput maximization for UAV-enabled
mobile relaying systems. \textit{IEEE Transactions on Communications, 64}(12), 4983--4996.

\noindent Zhou, Y., Barnes, C., Lu, J., Yang, J., \& Li, H. (2019). On the continuity of rotation
representations in neural networks. In \textit{Proceedings of the IEEE/CVF Conference on Computer
Vision and Pattern Recognition} (pp. 5745--5753).

\end{document}
